% chuong3.tex — sửa nội dung chương ở đây
\section{Phân tích vi phạm}\label{sec:3}

\subsection{Các hành vi thuộc nội dung khởi tố}\label{sec:3.1}

Trước hết, xét nhóm ba hành vi trực tiếp cấu thành mặt khách quan của tội phạm trong quyết định khởi tố của cơ quan điều tra, mỗi hành vi đều cho thấy sự xâm phạm đồng thời quy định hình sự và pháp luật về an toàn thông tin, bảo vệ dữ liệu cá nhân:

\subsubsection{Mua bán dữ liệu cá nhân trái phép}\label{sec:3.1.1}

\begin{itemize}
\item Hành vi thực tế: Từ cuối năm 2021 đến tháng 2/2026, sàn taphoammo.net cho phép hàng chục nghìn gian hàng công khai đăng tải, mua bán trái phép dữ liệu cá nhân (họ tên, số điện thoại, thông tin tài khoản ngân hàng, ví điện tử, thuế, bảo hiểm...) mà không có sự đồng ý của chủ thể dữ liệu.
\item \textit{Quy định bị vi phạm:} Vi phạm \textbf{khoản 5 Điều 7 Luật An toàn thông tin mạng 2015} (cấm kinh doanh trái pháp luật thông tin cá nhân – áp dụng cả giai đoạn cuối 2021 đến 02/2026), vi phạm \textbf{khoản 4 Điều 3} \textbf{Nghị định 13/2023/NĐ-CP} (\textit{“Dữ liệu cá nhân không được mua, bán dưới mọi hình thức”} – áp dụng từ 01/7/2023), và cấu thành tội phạm tại \textbf{điểm b khoản 1 Điều 288 Bộ luật Hình sự 2015 (sửa đổi 2017).}
\item \textit{Nhận xét:} Hành vi lập “chợ” dữ liệu của nhóm Đào Quang Huy đã bất hợp pháp ngay từ năm 2021 theo Luật An toàn thông tin mạng 2015 chứ không phải đợi đến khi Nghị định 13/2023/NĐ-CP có hiệu lực. Việc thương mại hóa dữ liệu cá nhân quy mô lớn đã trực tiếp cung cấp nguyên liệu đầu vào cho tội phạm giả mạo danh tính và lừa đảo chiếm đoạt tài sản.
\end{itemize}
\subsubsection{Mua bán tài khoản mạng xã hội và thư điện tử}\label{sec:3.1.2}

\begin{itemize}
\item \textit{Hành vi thực tế :} Sàn tổ chức cho các gian hàng mua bán hàng loạt tài khoản mạng xã hội (Facebook, TikTok, Telegram, Instagram) và thư điện tử (email); cơ quan công an đã phát hiện trường hợp đối tượng mua tài khoản trực tiếp từ sàn này để thực hiện hành vi lừa đảo.
\item \textit{Quy định bị vi phạm:} Vi phạm \textbf{khoản 1 (điểm b, c) Điều 288 Bộ luật Hình sự} (mua bán, trao đổi thông tin riêng hợp pháp trên mạng máy tính, mạng viễn thông mà không được phép của chủ sở hữu), đồng thời từ ngày 01/7/2023 vi phạm \textbf{khoản 4 Điều 3 Nghị định 13/2023/NĐ-CP}, bởi thông tin về tài khoản số của cá nhân thuộc danh mục \textbf{dữ liệu cá nhân cơ bản} theo \textbf{điểm k khoản 3 Điều 2 Nghị định 13/2023/NĐ-CP}.
\item \textit{Nhận xét:} Dù giới MMO thường ngụy trang đây là giao dịch “tài nguyên số” hay “tài khoản rác”, dưới góc độ pháp lý, thông tin định danh tài khoản số gắn với cá nhân là dữ liệu cá nhân cơ bản được pháp luật bảo vệ; việc mua bán các tài khoản này trực tiếp tiếp tay cho tội phạm ẩn danh trên không gian mạng.
\end{itemize}
\subsubsection{Cung cấp dịch vụ xác thực tài khoản hộ}\label{sec:3.1.3}

\begin{itemize}
\item \textit{Hành vi thực tế:} Bản tin của Bộ Công an (14/7/2026)~\cite{bca-2026} xác nhận taphoammo.net cho phép mở gian hàng quảng bá và giao dịch \textbf{“dịch vụ xác thực tài khoản”} bên cạnh dữ liệu cá nhân và tài khoản mạng xã hội.
\item \textit{Quy định bị vi phạm:} Đối chiếu trực tiếp với \textbf{Điều 288 Bộ luật Hình sự 2015 (sửa đổi 2017)} (\textit{Tội đưa hoặc sử dụng trái phép thông tin mạng máy tính, mạng viễn thông}).
\item \textit{Nhận xét:} Do cơ quan điều tra và báo chí chỉ nêu tên danh mục “dịch vụ xác thực tài khoản” mà không mô tả chi tiết phương thức kỹ thuật (như cho thuê SIM nhận OTP, cung cấp giấy tờ giả vượt eKYC hay can thiệp sinh trắc học), nhóm chỉ đối chiếu hành vi này với Điều 288 Bộ luật Hình sự để không suy luận vượt quá dữ kiện đã công bố.
\end{itemize}
Bên cạnh ba hành vi cốt lõi đã bị cơ quan điều tra trực tiếp khởi tố hình sự theo Điều 288 Bộ luật Hình sự nêu trên, quá trình rà soát cơ chế vận hành của taphoammo.net dưới góc độ quản lý nhà nước cho thấy nền tảng này còn bộc lộ nhiều dấu hiệu vi phạm các quy định pháp luật chuyên ngành khác.

\subsection{Các hành vi có dấu hiệu vi phạm khác}\label{sec:3.2}

Dựa trên các dữ kiện về kiểm duyệt gian hàng, hạ tầng máy chủ, tiếp thị đa nền tảng và thanh toán trung gian đã được Bộ Công an công bố, nhóm nghiên cứu xác lập thêm bốn hành vi có dấu hiệu vi phạm pháp luật chuyên ngành như sau:

\subsubsection{Kiểm duyệt người bán nhưng vẫn cho phép hàng hóa vi phạm lưu hành}\label{sec:3.2.1}

\begin{itemize}
\item \textbf{\textit{Dữ kiện:} Bộ Công an xác nhận sàn taphoammo.net có cơ chế kiểm duyệt người bán, bảo đảm giao dịch, giải quyết tranh chấp và thu phí 3\%–5\% mỗi giao dịch.}
\item \textit{Quy định liên quan:} Dữ kiện trên cho thấy chủ sàn có điều kiện biết bản chất hàng hóa bị cấm nhưng vẫn duy trì để hưởng lợi; hành vi này có dấu hiệu vi phạm \textbf{điểm a khoản 1 Điều 4} và nghĩa vụ của thương nhân, tổ chức cung cấp dịch vụ sàn giao dịch thương mại điện tử tại \textbf{khoản 8, 9, 11 Điều 36 Nghị định 52/2013/NĐ-CP (được sửa đổi, bổ sung bởi khoản 16 Điều 1 Nghị định 85/2021/NĐ-CP)} về trách nhiệm kiểm soát người bán, ngăn chặn và loại bỏ khỏi website thông tin mua bán hàng hóa, dịch vụ thuộc danh mục cấm kinh doanh.
\end{itemize}
\subsubsection{Máy chủ đặt ở nước ngoài, dấu hiệu không lưu trữ dữ liệu người dùng trong nước}\label{sec:3.2.2}

\begin{itemize}
\item \textit{Dữ kiện:} Nhóm vận hành dưới pháp nhân trong nước (Công ty HCT) nhưng thuê tên miền và đặt máy chủ tại nước ngoài để phục vụ hơn 1,35 triệu tài khoản.
\item \textit{Quy định liên quan:} Đối chiếu \textbf{khoản 3 Điều 26 Luật An ninh mạng 2018} và \textbf{Điều 26, Điều 27 Nghị định 53/2022/NĐ-CP}, doanh nghiệp trong nước phải tự động lưu trữ dữ liệu cá nhân, dữ liệu mối quan hệ và dữ liệu do người dùng tại Việt Nam tạo ra ở trong nước (tối thiểu 24 tháng; nhật ký hệ thống tối thiểu 12 tháng) mà không cần chờ văn bản yêu cầu như doanh nghiệp nước ngoài.
\end{itemize}
\subsubsection{Quảng cáo dịch vụ bị cấm trên Facebook, Telegram, TikTok, Instagram, Google}\label{sec:3.2.3}

\begin{itemize}
\item \textit{Dữ kiện:} Bộ phận marketing của sàn tổ chức quảng bá hình ảnh, dịch vụ trên Facebook, Telegram, TikTok, Instagram và Google.
\item \textit{Quy định liên quan:} Đối chiếu bản có hiệu lực tại thời điểm hành vi của \textbf{Luật Quảng cáo 2012} (đã sửa đổi bởi Luật số 75/2025/QH15), gồm \textbf{khoản 1 Điều 7} (hàng hóa, dịch vụ cấm kinh doanh thuộc diện cấm quảng cáo) và \textbf{khoản 1 Điều 8} (cấm quảng cáo các sản phẩm, dịch vụ quy định tại Điều 7). Vì việc mua bán dữ liệu cá nhân bị Luật ATTTM 2015 và Nghị định 13/2023/NĐ-CP cấm kinh doanh, nên việc chạy quảng cáo cho dịch vụ này có dấu hiệu vi phạm điều cấm của Luật Quảng cáo.
\end{itemize}
\subsubsection{Làm trung gian thanh toán cho giao dịch hàng hóa bị cấm, thu phí 3\% – 5\% mỗi giao dịch}\label{sec:3.2.4}

\begin{itemize}
\item \textit{Dữ kiện:} Sàn trực tiếp làm trung gian thanh toán, giữ tiền bảo đảm giao dịch và thu phí hoa hồng 3\%–5\% trên hơn 53 triệu giao dịch.
\item \textit{Quy định liên quan:} Đối chiếu \textbf{Điều 8 Nghị định 52/2024/NĐ-CP} (hiệu lực từ 01/7/2024) và văn bản tiền nhiệm giai đoạn trước 01/7/2024 là \textbf{Điều 6 Nghị định 101/2012/NĐ-CP (sửa đổi bởi Nghị định 80/2016/NĐ-CP)}. Nếu xét trường hợp sàn \textbf{không có giấy phép} của Ngân hàng Nhà nước thì vi phạm \textbf{khoản 7 Điều 8 Nghị định 52/2024/NĐ-CP} (cung ứng dịch vụ trung gian thanh toán không phép); và ngay cả với trường hợp khi \textbf{có giấy phép} vẫn vi phạm \textbf{khoản 8 Điều 8} (sử dụng hoặc tạo điều kiện sử dụng dịch vụ thanh toán/trung gian thanh toán để kinh doanh trái pháp luật) và có thể cả \textbf{khoản 12 Điều 8} (hoạt động không đúng nội dung giấy phép).
\end{itemize}
Mặc dù báo chí chưa công bố thông tin về việc đăng ký sàn thương mại điện tử với Bộ Công Thương, giấy phép trung gian thanh toán, hay các giao dịch phát sinh từ ngày 01/01/2026 đến tháng 02/2026 — yếu tố có thể khiến sàn chịu thêm mức phạt tối đa gấp 10 lần khoản thu từ việc mua bán dữ liệu theo khoản 6 Điều 7 và khoản 3 Điều 8 Luật Bảo vệ dữ liệu cá nhân 2025 (số 91/2025/QH15) — nhưng cả ba điểm khuyết thông tin này đều không làm ảnh hưởng đến việc xác lập bảy hành vi vi phạm đã nêu.

\subsection{Tổng hợp các vi phạm}\label{sec:3.3}

Từ việc đối chiếu chi tiết cả nhóm hành vi đã bị khởi tố hình sự lẫn nhóm hành vi có dấu hiệu vi phạm chuyên ngành nêu trên, toàn bộ kết quả định danh pháp lý của các lỗi vi phạm được đúc kết thành ma trận tổng hợp bên dưới nhằm tạo tiền đề phân tích các thách thức thực thi:

{\small\setstretch{1.1}
\begin{xltabular}{\textwidth}{>{\hsize=0.681\hsize\raggedright\arraybackslash}X>{\hsize=1.412\hsize\raggedright\arraybackslash}X>{\hsize=0.907\hsize\raggedright\arraybackslash}X}
\caption{Tổng hợp các hành vi vi phạm}\label{tab:3.5}\\
\toprule
\textbf{Hành vi} & \textbf{Quy định} & \textbf{Mức độ chứng thực} \\
\midrule
\endfirsthead
\toprule
\textbf{Hành vi} & \textbf{Quy định} & \textbf{Mức độ chứng thực} \\
\midrule
\endhead
\bottomrule
\endlastfoot
\textbf{Mua bán dữ liệu cá nhân trái phép} & \textbullet~Điểm b khoản 1 \& khoản 2 Điều 288 BLHS 2015 (sửa đổi 2017).\newline \textbullet~Khoản 5 Điều 7 Luật An toàn thông tin mạng 2015.\newline \textbullet~Khoản 4 Điều 3 Nghị định 13/2023/NĐ-CP (từ 01/7/2023).\newline \textbullet~Khoản 6 Điều 7 \& khoản 3 Điều 8 Luật BVDLCN 2025. & \textbf{Đã khởi tố hình sự} theo Điều 288 BLHS. \\
\textbf{Mua bán tài khoản mạng xã hội và thư điện tử} & \textbullet~Khoản 1 (điểm b, c) \& khoản 2 Điều 288 BLHS 2015 (sửa đổi 2017).\newline \textbullet~Khoản 4 Điều 3 kết hợp điểm k khoản 3 Điều 2 Nghị định 13/2023/NĐ-CP (từ 01/7/2023). & \textbf{Đã khởi tố hình sự} theo Điều 288 BLHS. \\
\textbf{Cung cấp dịch vụ xác thực tài khoản hộ} & \textbullet~Khoản 1 \& khoản 2 Điều 288 BLHS 2015 (sửa đổi 2017). & \textbf{Đã khởi tố hình sự} trong tội danh chung Điều 288 BLHS. \\
\textbf{Kiểm duyệt người bán nhưng vẫn cho phép hàng hóa vi phạm} & \textbullet~Điểm a khoản 1 Điều 4 \& khoản 8, 9, 11 Điều 36 Nghị định 52/2013/NĐ-CP (sửa đổi bởi khoản 16 Điều 1 Nghị định 85/2021/NĐ-CP). & Dữ kiện kiểm duyệt người bán, bảo đảm giao dịch và thu phí 3\%–5\% đã được công bố. \\
\textbf{Máy chủ đặt ở nước ngoài, không lưu dữ liệu người dùng trong nước} & \textbullet~Khoản 3 Điều 26 Luật An ninh mạng 2018.\newline \textbullet~Khoản 1, 2 Điều 26 \& khoản 1, 3 Điều 27 Nghị định 53/2022/NĐ-CP. & Đặt máy chủ ở nước ngoài là dữ kiện đã công bố; việc không lưu trữ dữ liệu người dùng tại Việt Nam là phân tích dựa trên hành vi vi phạm. \\
\textbf{Quảng cáo dịch vụ bị cấm trên Facebook, Telegram, TikTok, Instagram, Google} & \textbullet~Khoản 1 Điều 7 \& khoản 1 (kèm khoản 9) Điều 8 Luật Quảng cáo 2012 (bản có hiệu lực tại thời điểm hành vi). & Dữ kiện quảng bá trên 5 nền tảng đã được công bố; liên kết từ lệnh cấm kinh doanh dữ liệu cá nhân sang cấm quảng cáo. \\
\textbf{Làm trung gian thanh toán cho giao dịch hàng hóa bị cấm, thu phí 3\% – 5\%} & \textbullet~Khoản 7, 8, 12 Điều 8 Nghị định 52/2024/NĐ-CP (từ 01/7/2024).\newline \textbullet~Khoản 5, 6 Điều 6 Nghị định 101/2012/NĐ-CP sửa đổi bởi NĐ 80/2016/NĐ-CP (trước 01/7/2024). & Dù có hay không có giấy phép của NHNN thì vẫn cấu thành hành vi vi phạm. \\
\end{xltabular}}

Như vậy, từ bảng tổng hợp, ta có thể giải quyết Câu hỏi nghiên cứu số 1 bằng cách chỉ rõ 7 hành vi vi phạm cụ thể, đồng thời phác họa rõ những thách thức về bảo mật và kỹ thuật, pháp lý và quản trị liên ngành, tạo cơ sở đầu vào trực tiếp để Chương 4 tiếp tục phân tích nguyên nhân khiến sàn tồn tại kéo dài trước khi bị triệt phá.

